% =========================================================================
% ปฏิบัติการที่ 1: สถาปัตยกรรม WebAR และเครื่องมือพัฒนา
% =========================================================================

\chapter{สถาปัตยกรรม WebAR และเครื่องมือพัฒนา}
\label{ch:lab01}

\noindent{\large\bfseries\color{cyberblue} การวางโครงสร้าง Three.js, A-Frame และ WebXR Device API บนเว็บเบราว์เซอร์}

\vspace{0.4cm}

\begin{labobjectivebox}{สถาปัตยกรรม WebAR และเครื่องมือพัฒนา}
\textbf{วัตถุประสงค์การเรียนรู้}
\begin{enumerate}[leftmargin=1.8cm, itemsep=2pt]
    \item เข้าใจสถาปัตยกรรมระบบความจริงเสริมบนเว็บ (Web-based Augmented Reality Architecture)
    \item สามารถติดตั้งและกำหนดค่า Three.js และ A-Frame เพื่อเรนเดอร์ฉาก 3 มิติ
    \item สามารถควบคุมวงรอบการเรนเดอร์กราฟิกให้มีอัตราการแสดงผลระดับ 60 เฟรมต่อวินาที
\end{enumerate}

\vspace{4pt}
\textbf{เครื่องมือและอุปกรณ์จำเป็น} \enspace เว็บเบราว์เซอร์สมัยใหม่ (Chrome/Edge), Visual Studio Code, Three.js r128, A-Frame v1.4.0
\end{labobjectivebox}

\section{หลักการและทฤษฎีพื้นฐาน}

เทคโนโลยีความจริงเสริมบนเว็บเบราว์เซอร์ (Web-based Augmented Reality หรือ WebAR) เป็นวิวัฒนาการสำคัญที่ช่วยให้ผู้ใช้งานสามารถเข้าถึงประสบการณ์เสมือนจริงได้ทันทีผ่านโปรแกรมเปิดเว็บโดยไม่ต้องดาวน์โหลดและติดตั้งแอปพลิเคชันที่มีขนาดใหญ่ สถาปัตยกรรมของ WebAR อาศัยการทำงานร่วมกันของมาตรฐานเปิดระดับสากล ได้แก่ WebGL สำหรับการประมวลผลกราฟิกฮาร์ดแวร์เร่งความเร็ว และ WebXR Device API ซึ่งทำหน้าที่ตรวจจับและเชื่อมต่อข้อมูลเชิงพื้นที่จากเซนเซอร์ของอุปกรณ์เข้ากับฉากทัศน์เสมือน

หัวใจหลักของกราฟิก 3 มิติบนเว็บคือเอนจิน Three.js ซึ่งทำหน้าที่เป็นตัวประสานงานระดับสูง (High-level Abstraction) ครอบ WebGL ช่วยลดความซับซ้อนในการเขียนภาษาเฉดเดอร์ (GLSL) โดยแบ่งโครงสร้างออกเป็น 3 ส่วนหลัก ได้แก่ ฉากทัศน์ (Scene) ทำหน้าที่เป็นพื้นที่รองรับวัตถุและโมเดล กล้องเสมือน (Camera) เช่น กล้องมุมมองแบบเพอร์สเปกทีฟ (Perspective Camera) ที่จำลองสายตาของมนุษย์ และตัวเรนเดอร์ (WebGLRenderer) ทำหน้าที่คำนวณและวาดภาพลงบนผืนผ้าใบแบบเฟรมต่อเฟรม ผ่านฟังก์ชัน \texttt{requestAnimationFrame()}

\begin{conceptbox}{สาระสำคัญของปฏิบัติการที่ 1}
การเข้าใจโครงสร้างทางคณิตศาสตร์และการประมวลผลเชิงพื้นที่ ช่วยให้นักศึกษาสามารถออกแบบระบบเสมือนจริงที่ตอบสนองต่อผู้ใช้งานได้อย่างเป็นธรรมชาติและแม่นยำสูง ทั้งยังสามารถต่อยอดสู่การพัฒนาแอปพลิเคชันทางวิทยาศาสตร์และอุตสาหกรรมได้อย่างมีประสิทธิภาพ
\end{conceptbox}

\section{ขั้นตอนการปฏิบัติการและการทดลอง}

ให้นักศึกษาดำเนินการสร้างไฟล์โครงการและเขียนคำสั่งตามลำดับขั้นตอนต่อไปนี้ โดยตรวจสอบความถูกต้องของไลบรารีที่เรียกใช้งานและเส้นทางไฟล์

\begin{codebox}{โครงสร้างโค้ดหลักของปฏิบัติการที่ 1}
\begin{lstlisting}[language=HTML]
<!DOCTYPE html>
<html lang="th">
<head>
    <meta charset="UTF-8">
    <title>WebAR Lab 01 - Three.js Fundamentals</title>
    <script src="https://cdnjs.cloudflare.com/ajax/libs/three.js/r128/three.min.js"></script>
    <style>
        body { margin: 0; overflow: hidden; background: #030712; }
        canvas { width: 100vw; height: 100vh; display: block; }
    </style>
</head>
<body>
    <div id="webgl-container"></div>
    <script>
        // 1. กำหนดตัวแปรหลักสำหรับฉากทัศน์ กล้อง และตัวเรนเดอร์
        const scene = new THREE.Scene();
        const camera = new THREE.PerspectiveCamera(75, window.innerWidth / window.innerHeight, 0.1, 1000);
        const renderer = new THREE.WebGLRenderer({ antialias: true, alpha: true });
        renderer.setSize(window.innerWidth, window.innerHeight);
        renderer.setPixelRatio(window.devicePixelRatio);
        document.getElementById('webgl-container').appendChild(renderer.domElement);

        // 2. สร้างวัตถุเรขาคณิต 3 มิติแบบโฮโลกราฟิก
        const geometry = new THREE.BoxGeometry(1.2, 1.2, 1.2);
        const material = new THREE.MeshStandardMaterial({
            color: 0x06b6d4,
            wireframe: false,
            metalness: 0.8,
            roughness: 0.2
        });
        const cube = new THREE.Mesh(geometry, material);
        scene.add(cube);

        // 3. กำหนดแหล่งกำเนิดแสงในฉาก
        const ambientLight = new THREE.AmbientLight(0xffffff, 0.7);
        scene.add(ambientLight);
        const pointLight = new THREE.PointLight(0x38bdf8, 2, 50);
        pointLight.position.set(5, 5, 5);
        scene.add(pointLight);

        camera.position.z = 3.5;

        // 4. วงรอบการเรนเดอร์แบบ 60 FPS
        function animate() {
            requestAnimationFrame(animate);
            cube.rotation.x += 0.01;
            cube.rotation.y += 0.015;
            renderer.render(scene, camera);
        }
        animate();
    </script>
</body>
</html>
\end{lstlisting}
\end{codebox}

\section{การทดสอบระบบและการบันทึกผล}

เมื่อรันโปรแกรมบนเซิร์ฟเวอร์จำลอง (Local Development Server) ให้ทำการทดสอบการทำงานตามเกณฑ์ประเมินในตารางต่อไปนี้ พร้อมบันทึกผลการสังเกตการณ์ลงในรายงานผลการทดลอง

\begin{table}[htbp]
\centering
\small
\begin{tabularx}{\textwidth}{c X c Y}
\toprule
\textbf{ลำดับ} & \textbf{รายการทดสอบและพฤติกรรมที่คาดหวัง} & \textbf{เกณฑ์ผ่าน} & \textbf{ผลการทดสอบ} \\
\midrule
1 & การเริ่มต้นระบบและการเข้าถึงสิทธิ์ฮาร์ดแวร์ & ภายใน 3 วินาที & ผ่าน / ไม่ผ่าน \\
2 & ความเสถียรของการตรวจจับและการคำนวณพิกัด & อัตราความแม่นยำ $> 90\%$ & ผ่าน / ไม่ผ่าน \\
3 & อัตราการแสดงผลกราฟิกต่อเนื่อง (Framerate) & ไม่ต่ำกว่า 55 FPS & ผ่าน / ไม่ผ่าน \\
4 & การตอบสนองต่อปฏิสัมพันธ์เชิงพื้นที่ & ความหน่วง $< 40$ ms & ผ่าน / ไม่ผ่าน \\
\bottomrule
\end{tabularx}
\caption{ตารางประเมินผลการทำงานของระบบในปฏิบัติการที่ 1}
\label{tab:eval_lab01}
\end{table}

\section{ภารกิจท้าทายและการต่อยอด}

\begin{activitybox}{กิจกรรมส่งเสริมทักษะขั้นสูง}
ให้นักศึกษาปรับแต่งขนาดรูปทรงเรขาคณิตให้เป็นทรงกลม (SphereGeometry) และเพิ่มแหล่งกำเนิดแสงสีส้มทอง (Amber Light) อีก 1 จุด พร้อมคำนวณอัตราการใช้หน่วยความจำและการตอบสนองต่อการปรับขนาดหน้าจอแบบ Responsive
\end{activitybox}

\section{คำถามท้ายการทดลอง}

ให้นักศึกษาตอบคำถามเชิงวิเคราะห์ต่อไปนี้ลงในสมุดบันทึกปฏิบัติการ

\begin{enumerate}[leftmargin=1.8cm, itemsep=6pt]
    \item สถาปัตยกรรมของ WebXR Device API แตกต่างจากการเรียกใช้ WebGL ดั้งเดิมอย่างไรในแง่ของการจัดการพิกัดเชิงพื้นที่
    \item เหตุใดการตั้งค่า `alpha: true` ใน WebGLRenderer จึงมีความจำเป็นอย่างยิ่งสำหรับการพัฒนาแอปพลิเคชันแบบความจริงเสริม
    \item จงอธิบายบทบาทของฟังก์ชัน requestAnimationFrame() ต่อการรักษาอัตราเฟรมเรตและการประหยัดพลังงานประมวลผล
\end{enumerate}

